\chapter{Introduction}

\section{Background and Motivation}

Portland cement is the most commonly used binder for the production of concrete.
After adding water to the cement, a reaction called `hydration' begins.
The cement reactivity is determined by its composition and particle size.
During hydration, the cement begins to dissolve and new crystals start to form (see Figure~\ref{fig:intro}).
This ultimately causes hardening of the cement paste.
Furthermore, this microstructure determines the properties of the resulting concrete.

In 1929 \textsc{Anderegg} and \textsc{Hubbell} introduced their paper as follows:
\begin{iquote}
	The lack of knowledge concerning the rate of hydration of cement clinker has hindered a more complete understanding of the complicated setting and hardening reactions that cements undergo. \cite{anderegg.1929}
\end{iquote}


\begin{figure*}[b!]
    \centering

	\ifdefined\PRINTABLE
		\includegraphics[width=\textwidth]{intro-anim/Animation06.jpg}
	\else
		\animategraphics[palindrome,autoplay,width=\textwidth]{6}{intro-anim/Animation}{01}{10}
	\fi
    \caption{Artistic 3D reconstruction of \acl*{CSH}-fibres within a capillary pore of \SI{28}{\day} hydrated \acl*{C3S}. The dataset was recorded using a \acl*{SEM} combined with a \acl*{FIB}. \cite{kleiner.2021c}}
    \label{fig:intro}
\end{figure*}

While the knowledge and measurement techniques have significantly improved within the last century, this is still true to a certain extent today.
Good knowledge of the structural development and therefore, reliable measurement and evaluation methods are crucial to understand, model \cite{nguyen-tuan.2023,nguyen-tuan.2024} and consequently influence the hydration of cements.
Good measurement and evaluation methods are therefore required.

A \ac{SEM} is a powerful tool to image the microstructure of a lot of materials.
Electron microscopy is widely known for its \ac{SE}-mode, which gives a great idea of the surface profile of the specimen, or the \ac{BSE}-mode resulting in a chemical contrast image.
Nevertheless, the data extracted from the resulting images can be greatly improved, if also data from \ac{EDX}, \ac{EBSD}, optical microscopy, \ac{mCT}, laser ablation or nano indentation are recorded.
The combination of multiple (imaging) techniques as the afore mentioned is often referred to as correlative microscopy.
It results in significantly more information than the simple sum of the individual method would be able to produce.

Therefore, this thesis is devoted to improve the characterization the micro-structure of unhydrated and hydrated cement based binders by combining a set of imaging techniques like \ac{SEM} and \ac{mCT} and comparing them to selected methods like \ac{LTDSC} and \ac{MIP}.
The goals are to better understand the influence of the devices and techniques, and finding improved ways to combine them to gain better knowledge about unhydrated and hydrated cements and the cement hydration, defining protocols and best practices to obtain statistical relevant data and to develop or adopt software tools for image analysis.
Furthermore, the composition of clinker and the structure and porosity of hydration products were the focus of this work.

Modern computer vision methods allow to develop complicated and creative evaluation methods, which were computational to expensive, just a few years ago.
Therefore, existing analytical methods are to be improved and further developed to be reproducible and and fairly easily repeatable.

Please note, that this work uses the cement chemist notation for most chemical equations (see Table~\ref{tab:ccn}).

% TODO: Christane:
% hier könnte schon noch so ein kleiner abriß stehen, dass man mit mikroskopie schon immer versucht materialien wie beton besser zu verstehen und eigenschaften verständlich bzw vorhersagbar zu machen. dass man dank moderner mikroskopie nun modelle entwickelnn kann um die mikrostrutur zu simulieren (sheet growth chapter hierher holen?, hymostruc oder eben auch 2D-3D transformation in zukunft dank ML blabla nutzen kann..)

{\color{red}
Chapter ideas
\begin{itemize}
	
	\item Dissolution- and hydration rates of clinker phases depending on the crystal orientation (SEM+EBSD+VSI)
	\item Chemical and structural 3D-analysis of hydrated cementitious materials
\end{itemize}
}

% Einleitung/Grundlagen:
% alles mit ausreichend Literatur belegen, sehr allgemein:
% Was ist Zement? Verwendung? Gesellschaftlicher Nutzen? Zeitlicher Kontext?

% spezieller Klinkerphasen/Bindemittel:
% Was sind Klinkerphasen? Wozu Binder? Alternativen?
% Wie werden beide methodisch untersucht?
% Unterschiede? Herausforderungen? Probenbehandlung?
% historische und moderne Methodik?
% 2D- und 3D-Methoden? Unterschiede? Vorteile?